\documentclass[11pt]{article}
\usepackage[margin=2.2cm]{geometry}
\usepackage[T1]{fontenc}
\usepackage{amsmath,amssymb,booktabs,graphicx,xcolor,enumitem,caption,array}
\usepackage[section]{placeins}
\usepackage{hyperref}
\setlist{itemsep=2pt,topsep=3pt}
\hypersetup{colorlinks=true,linkcolor=blue!50!black,urlcolor=blue!50!black}
\captionsetup{font=small,labelfont=bf}
\setlength{\emergencystretch}{3em}
\newcommand{\fig}[2]{\begin{figure}[ht]\centering\includegraphics[width=\textwidth]{figures/#1}\caption{#2}\end{figure}}
\newcommand{\read}[1]{\par\noindent\textbf{How to read it:} #1\par}
\newcommand{\shows}[1]{\par\noindent\textbf{What it shows:} #1\par\medskip}

\title{How well do the level-1 models work? A plain-language guide to the v2 results}
\author{ccd\_opt\_ml}
\date{2026-11-10}

\begin{document}
\maketitle

\section*{In one page}
We image a CCD at many settings and measure, from each image, its \textbf{noise} ($\sigma_e$), its \textbf{dark
charge} ($\lambda$) and its \textbf{gain}. Images are slow (about 10 minutes), so we tried to predict them from things
that are fast: the \emph{settings} we chose, and a short \emph{trace} (0.2\,s of the output signal). The findings:
\begin{enumerate}
\item \textbf{The dark charge is mostly set by two things: which amplifier it is, and how long the readout takes
  (\texttt{delayInt}).} Four amplifiers of the same CCD, at the same settings, have dark charges from 0.20 to 0.50
  e$^-$/pixel. On amplifiers they know, the models predict the dark charge mainly by recognising the amplifier (a trace
  gives it away 99\,\% of the time) and reading \texttt{delayInt}.
\item \textbf{That is why they fail on an amplifier they have never seen}: they cannot know its level, so they guess
  roughly the average of the others. A low-dark-charge amplifier is then predicted 45\,\% too high, a high one 50\,\% too
  low. \textbf{One image of the new amplifier fixes this} (error 48--53\,\% $\to$ 17--18\,\%).
\item \textbf{Predicting a \emph{change} avoids the problem}: the level cancels when you compare two settings of the
  same amplifier. The change models get the direction of a dark-charge change right 97--98\,\% of the time on a new
  amplifier, as on a known one.
\item \textbf{The integration time \texttt{delayInt} does most of the work}: alone it explains 76\,\% of the dark-charge
  spread and 51\,\% of the noise spread. Without it, models that use only the settings fall to 50--54\,\% (a coin flip)
  on the changes where it moved. When \texttt{delayInt} stays the same, the trace models still call the noise direction
  71--93\,\% of the time: the trace sees the noise itself.
\item \textbf{We never had data to train on ``timing fixed, voltages changed''}: in v2 every change of voltage came with
  a change of timing. The v3 noise set is the first data of that kind.
\end{enumerate}

\section{The words used}
\begin{description}
\item[Image] The slow, full measurement (about 10 minutes). It gives the true values (``labels''): noise $\sigma_e$
  (electrons; lower is better), dark charge $\lambda$ (electrons per pixel; lower is better), gain (ADU per electron).
\item[Trace] A fast recording (0.2\,s) of the output signal. It contains the noise directly, but almost no dark charge.
\item[Settings] What we set: timing (\texttt{delayInt} = the integration time of each sample, and other delays) and
  voltages (clock levels, VSUB, etc.). In v2 all of them were changed together in every packet.
\item[Packet] One measurement at one setting: traces, a short strip and an image, on four amplifiers at once.
\item[Amplifier] Each CCD has four outputs (ch0--ch3). They see the same settings but behave a little differently.
\item[Absolute model] Predicts the value (e.g. $\lambda = 0.31$) at a setting.
\item[Change model] Predicts the \emph{change} between two packets of the same amplifier (e.g. ``$\lambda$ goes up by
  20\,\%''), or only its direction (up or down).
\item[Settings model / trace model] A model that sees only the settings / only the traces. Hybrid models see both.
\item[Classifier] A change model that answers only ``up or down?'', with a probability.
\item[New amplifier] The test that matters for a new CCD: the model is trained on the other amplifiers and tested on
  one it has never seen.
\item[Clear change] A change bigger than the measurement can resolve: $\sigma_e$ $>$ 5\,\%, $\lambda$ $>$ 15\,\%, gain
  $>$ 1.2\,\%. ``Direction right'' is counted on these only.
\item[Typical error] The root-mean-square error in \% (log units). 10\,\% means a prediction is typically 10\,\% off.
\end{description}

\section{How everything was tested}
All numbers are from session D of v2 (127 measurements on three amplifiers, ch0, ch1, ch3). The night is cut into five
blocks in time. A model is always trained on four blocks and tested on the fifth, and for the ``new amplifier'' test
also never sees the tested amplifier. So every prediction is for packets the model has not seen.

\section{Which amplifier is it? How the models predict the dark charge}

\fig{amp_levels.png}{Each amplifier at the same (nominal) settings.}
\read{each bar is one amplifier, all read out with identical settings; the line on top is the spread of repeated
measurements.}
\shows{the dark charge is 0.20 (ch0), 0.26 (ch1), 0.40 (ch3), 0.50 (ch2) e$^-$/pixel at the same settings: a factor 2.5.
The noise differs by about 20\,\%, the gain by 6\,\%. Each amplifier has its own \emph{level}.}

\fig{amp_share.png}{How much of each quantity is set by the amplifier, by the settings, or by both.}
\read{1 would mean ``explains everything''. Red: knowing only which amplifier; blue: knowing only the settings; green:
both.}
\shows{for the dark charge, the amplifier alone explains 50\,\%, the settings alone 34\,\%, both together 88\,\%. For the
gain the amplifier is almost everything (76\,\%). For the noise the settings matter more (40\,\%) than the amplifier
(18\,\%).}

\fig{amp_from_trace.png}{Can a model tell which amplifier a trace came from?}
\read{a classifier sees one trace and names the amplifier; tested on packets it never saw. Guessing would give 25\,\%.}
\shows{99\,\% right, even after removing the trace features we already knew were amplifier ``fingerprints''. So a trace
model always effectively knows which amplifier it is looking at.}

\fig{amp_known_vs_new.png}{Absolute models on amplifiers they know (green) and on a new one (red).}
\read{lower is better. The dashed line: a model told the settings and which amplifier it is (no trace).}
\shows{on known amplifiers, the trace models predict the dark charge to 13--23\,\%. On a new amplifier they fall to
50--83\,\%. A model that is simply \emph{told} the amplifier and the settings gets 5\,\%. So on known amplifiers the dark
charge is predicted mostly by recognising the amplifier and reading the settings, not by measuring dark charge in the
trace (a trace hardly contains any).}

\begin{center}
\begin{tabular}{lrrr}
\toprule
dark charge, new amplifier: average error of the prediction & ch0 & ch1 & ch3\\
\midrule
settings model & $+45$\,\% & $+9$\,\% & $-53$\,\%\\
trace physics model & $+122$\,\% & $+55$\,\% & $-46$\,\%\\
trace trees & $+57$\,\% & $+39$\,\% & $-51$\,\%\\
\bottomrule
\end{tabular}
\end{center}
A new amplifier gets roughly the level of the others: ch0, the lowest, is predicted far too high; ch3, the highest, far
too low; ch1, in the middle, comes off best.

\fig{amp_calibration_image.png}{A new amplifier after 0, 1, 2, 4, 8 images of it (the model learns its level from them).}
\read{left of each panel: no image of the new amplifier; moving right: more of its images.}
\shows{with no image, the dark-charge error is 48\,\% (ch0) and 53\,\% (ch3) but only 16\,\% for ch1, whose level is near
the average. \textbf{One image brings all three to 17--18\,\%}; more images add little. The same for the gain (3.5--4.5\,\%
$\to$ 1.1--1.3\,\%). So the amplifier-to-amplifier problem is mostly a \emph{level}, and one image measures it.}

\fig{amp_changes.png}{Change models on amplifiers they know (green) and on a new one (red).}
\read{higher is better; ``direction right'' on clear changes between two packets of one amplifier.}
\shows{almost no loss on a new amplifier: dark charge 97--98\,\% (known 98--99\,\%), noise 88--90\,\% with the trace
models (settings only 69\,\%). Comparing two settings of the \emph{same} amplifier cancels its level. That is why the change
models carry over to new amplifiers and the absolute models do not.}

\section{The integration time (\texttt{delayInt})}

\fig{delayint_effect.png}{What \texttt{delayInt} does (top) and how much of the changes is left without it (bottom).}
\read{top: each dot is a measurement, relative to its amplifier's average; lines are smooth fits per amplifier. Bottom:
the spread of changes between two packets: all changes (blue), with the part \texttt{delayInt} explains removed (orange),
and pairs where \texttt{delayInt} stayed the same (green). Dotted lines: the ``clear change'' size.}
\shows{the noise falls with a longer integration time and flattens above about 1500 cycles; the dark charge rises
steadily (longer readout, more dark charge). \texttt{delayInt} alone explains 51\,\% of the noise spread and 76\,\% of the
dark-charge spread; nothing of the gain (the negative number means ``no better than the average''). Removing its part
shrinks the typical change from 27 to 16\,\% (noise) and from 40 to 17\,\% (dark charge): what is left is the other
settings and the measurement scatter. All amplifiers respond the same way: the lines lie on top of each other.}

\fig{delayint_accuracy.png}{How often the change models are right, against how much \texttt{delayInt} changed.}
\read{left on each axis: pairs where \texttt{delayInt} stayed the same (within 7\,\%); right: it doubled or more.
Dashed lines: settings-only models.}
\shows{when \texttt{delayInt} changes a lot, almost every model is right ($\geq$ 96\,\%). When it stays the same: for the
noise the settings-only models are near a coin flip (62--69\,\%), while the trace models get 69--85\,\% and the hybrid
(settings + traces) 87--93\,\%. For the dark charge, everything is right on the few pairs where \texttt{delayInt} stayed
the same (40 pairs; there \texttt{delayRGH}, another delay, did the work).}

\fig{delayint_removed.png}{Experiment: what happens when the integration time is taken away.}
\read{orange: test pairs where \texttt{delayInt} stayed the same; blue: test pairs where it changed. Three set-ups:
trained on all pairs; trained \emph{only} on pairs with \texttt{delayInt} unchanged; a settings model not told
\texttt{delayInt}.}
\shows{
\begin{itemize}
\item a settings model \emph{not told} \texttt{delayInt} falls to 50\,\% (noise) and 54\,\% (dark charge) on the changes
  where it moved: a coin flip. Most of what the settings models know is \texttt{delayInt}.
\item trained \emph{only} on pairs with \texttt{delayInt} unchanged (144 training pairs, almost all of them a setting
  against itself), the models learn little: on noise 53\,\% (settings) and 76\,\% (traces) on such pairs, against 58\,\% and
  71\,\% when trained on everything. v2 simply has no data of this kind to learn from.
\end{itemize}}

\section{Diagnostics of the change models}

\fig{change_scatter_sigma_e.png}{Every change model: predicted against true change of the noise (new amplifier).}
\read{each dot is a pair of packets. Perfect = on the diagonal. Dots in the top-right and bottom-left squares have the
direction right. Orange: \texttt{delayInt} unchanged. For the two classifiers the vertical axis is ``how sure it is that
the noise goes up'' (top = sure up, bottom = sure down).}
\shows{the settings-only models scatter widely (69--73\,\% direction, typical error 27--30\,\%). The trace models lie close
to the diagonal (85--91\,\%, error 16--18\,\%). The orange dots sit near the centre: with \texttt{delayInt} fixed the
changes are small and harder to call. The classifiers are usually certain (dots at the top and bottom edges) and
usually right. The same figures exist for the dark charge and the gain (\texttt{change\_scatter\_lambda.png},
\texttt{change\_scatter\_gain.png}).}

\fig{direction_vs_size.png}{How big a change must be before each model calls it right.}
\read{left = small changes, right = large ones. 50\,\% is a coin flip. Dotted vertical line: the ``clear change'' size.}
\shows{every model is near a coin flip (50--70\,\%) on changes smaller than the measurement scatter. For the dark charge
almost every model passes 90\,\% from changes of about 15--20\,\% and reaches 100\,\% above 40\,\%. For the noise the trace
models pass 90\,\% from changes of about 20\,\%, while the settings-only models (dashed) stay at 76--86\,\% even for the
largest changes: the traces resolve noise changes that the settings cannot explain.}

\fig{classifier_reliability.png}{The two direction classifiers: is their stated confidence right?}
\read{if the classifier says ``70\,\% sure it goes up'', it should be right 70\,\% of the time: points on the diagonal.
Grey bars: how many pairs got each confidence. ``ROC area'': 0.5 = useless, 1 = perfect.}
\shows{ROC area 0.95 (noise), 0.99--1.00 (dark charge), 0.92--0.94 (gain): they separate up from down very well. Most
pairs get a confident answer at the ends.}

\fig{size_calibration.png}{Do the size models get the \emph{size} of a change right?}
\read{each point averages the pairs with a similar true change. On the dotted line = right size. Flatter than the line =
the model predicts changes too small.}
\shows{every model predicts the changes too small. Slopes for the noise: 0.35--0.38 for the settings models (they
predict a third of the real change) and 0.58--0.70 for the trace models; for the dark charge 0.71--0.93. For ranking
settings this does not matter (the order is kept); for quoting an expected improvement it does.}

\fig{by_pair_type.png}{Direction by kind of pair.}
\read{within packet: the same packet's reference and new traces; $n\to n{+}1$, $n\to n{+}2$: neighbouring packets;
random: any two packets of the night.}
\shows{for the noise the trace models get 83--95\,\% on every kind of pair; the settings-only model 65--85\,\%, lowest
on $n\to n{+}2$ (65\,\%) and within-packet (67\,\%) pairs. For the dark charge every model is 92--100\,\% except on the
strip pairs (71--86\,\%: the strip's dark charge is a different, noisier measurement).}

\fig{by_amplifier.png}{Direction for each amplifier when it is the new one.}
\shows{no amplifier is much harder than the others: on the noise the trace models get 84--93\,\% whichever amplifier is
new (settings only 65--73\,\%); on the dark charge 96--100\,\%.}

\section{Diagnostics of the absolute models}

\fig{abs_pred_vs_true_lambda.png}{Absolute models: predicted against measured dark charge, new amplifier.}
\read{each colour is one amplifier; on the dotted line = perfect.}
\shows{each colour forms its own line, shifted from the diagonal: the shape is right, the level is wrong. That is the
amplifier level of Section 3. The noise and gain versions are \texttt{abs\_pred\_vs\_true\_sigma\_e.png} and
\texttt{\_gain.png}.}

\fig{abs_residuals.png}{Where the absolute models miss: against \texttt{delayInt} (top) and against time (bottom).}
\read{0 = no miss; lines are running averages.}
\shows{noise: every model predicts 5--20\,\% too low at the shortest integration times; at the longest the trace
models predict 20--25\,\% too high (the settings model 10\,\%). Over the night the noise misses wander by about $\pm$10\,\%.
Dark charge: the misses are dominated by the amplifier level (each model's line sits 0--50\,\% above zero over the whole
night, the new-amplifier offset of Section 3) rather than by \texttt{delayInt} or time. Gain: within a few \%.}

\section{What this means}
\begin{itemize}
\item For predicting values (absolute models), the hard part is not the settings but the \textbf{amplifier's level}.
  Recognising a known amplifier is easy (99\,\%); guessing the level of an unknown one is not. \textbf{One image of a new
  amplifier (or CCD) solves most of it.}
\item For comparing settings (change models), the level cancels: they carry over to new amplifiers.
\item \textbf{\texttt{delayInt} dominates v2.} The settings models are mostly \texttt{delayInt} models. The traces add the
  most where \texttt{delayInt} does not change, but v2 has few such pairs and none where only voltages changed.
\item \textbf{The next question} --- how the voltages move the noise and the dark charge at a fixed integration time ---
  needs the v3 noise set (voltage-only packets at nominal timing). v2 cannot answer it.
\end{itemize}

\section*{Files}
Figures and tables: \texttt{processed\_v2\_trainD/lvl1\_results/diagnostics/} (\texttt{diagnostics.md},
\texttt{amplifier\_identity.md}). Scripts: \texttt{lvl1/diagnostics\_v2.py}, \texttt{lvl1/amplifier\_identity.py}.

\end{document}
